\begin{summarybox}
	\textbf{\textcolor{CSummary}{一句话核心}}

	\textbf{椭圆焦半径公式将点到焦点的距离表示为横坐标的线性函数：$|PF_1|=a+ex_0,\ |PF_2|=a-ex_0$。}

	\medskip
	\textbf{\textcolor{CSummary}{使用条件}}
	\begin{itemize}[leftmargin=*, itemsep=0.5em, topsep=0.4em]
		\item \textbf{条件1（焦点位置）：} 椭圆焦点必须在 $x$ 轴上，标准方程为 $\frac{x^2}{a^2}+\frac{y^2}{b^2}=1\ (a>b>0)$。
		\item \textbf{条件2（点在椭圆）：} 点 $P(x_0,y_0)$ 必须在椭圆上，满足椭圆方程。
	\end{itemize}

	\medskip
	\textbf{\textcolor{CSummary}{关键提醒}}
	\begin{itemize}[leftmargin=*, itemsep=0.5em, topsep=0.4em]
		\item \textbf{易错点（记准符号）：} 左焦半径对应加号，右焦半径对应减号。
		\item \textbf{检查项（核实范围）：} 使用公式前检查 $x_0$ 是否在 $[-a,a]$ 内，防止公式误用。
	\end{itemize}


\medskip
\textbf{\textcolor{CSummary}{三条证明路径（优先记第二条）}}
\begin{itemize}[leftmargin=*, itemsep=0.35em, topsep=0.3em]
 \item \textbf{最快：} 第二定义 $\Rightarrow |PF_2|=e(\frac{a^2}{c}-x_0)=a-ex_0$，左焦点同理。
 \item \textbf{最低门槛：} $r_1+r_2=2a$，$r_1^2-r_2^2=4cx_0$，所以 $r_1-r_2=2ex_0$；两式联立即得结论。
 \item \textbf{常规：} 距离公式 $\Rightarrow$ 用椭圆方程消去 $y_0^2$ $\Rightarrow$ 配方开方。
\end{itemize}

\medskip
\textbf{\textcolor{CSummary}{考场触发器}}
\textbf{见到“两焦点距离＋点的横坐标”就考虑焦半径；见到“两焦点距离的和与差”就考虑平方差法。}
记号判断：点向右移，离左焦点相对更远，所以\textbf{左加右减}。
\[
 \boxed{|PF_1|-|PF_2|=2ex_0},\qquad
 \boxed{|PF_1|\,|PF_2|=a^2-e^2x_0^2\in[b^2,a^2]}.
\]
焦点在 $y$ 轴：换成\textbf{下加上减}，把 $x_0$ 换成 $y_0$。
\end{summarybox}
